% Einheit 1 – Kreisumfang – Nr. 10–21
\setcounter{aufgabe}{9}
\einheitenkopf[e1][1 Kreisumfang]{Einheit 1 von 3 · Kreisumfang}
\zweigzeile{Hier lernst du, den Umfang eines Kreises aus Radius oder Durchmesser zu berechnen und umgekehrt · neu oder schon bekannt – je nach Buch · P10 · baut auf: Kommazahlen und Runden, Umfang und Fläche, Formeln umstellen}

\begin{aufgabe}{Ich erkenne, ob Radius oder Durchmesser gegeben ist.}
Kreuze an: Ist die Strecke ein Radius, ein Durchmesser oder keins von beiden?

\begin{kreis}[1.8]
\mittelpunkt{M}
\radius{40}{1}
\durchmesser{150}{2}
\sehne{200}{280}{3}
\end{kreis}

\begin{teile}
\teil Strecke 1: \quad \kreuz{Radius}\kreuz{Durchmesser}\kreuz{keins von beiden}
\teil Strecke 2: \quad \kreuz{Radius}\kreuz{Durchmesser}\kreuz{keins von beiden}
\teil Strecke 3: \quad \kreuz{Radius}\kreuz{Durchmesser}\kreuz{keins von beiden}
\teil Lena misst ein rundes Planschbecken von Rand zu Rand, genau durch die Mitte. \\
\kreuz{Radius}\kreuz{Durchmesser}\kreuz{keins von beiden}
\teil Der Zirkel ist beim Zeichnen des Kreises 4 cm weit geöffnet. \\
\kreuz{Radius}\kreuz{Durchmesser}\kreuz{keins von beiden}
\end{teile}
\end{aufgabe}

\begin{aufgabe}{Ich erkenne, ob der Rand oder die Fläche gefragt ist.}
Kreuze an: Wird der Umfang (Rand) oder die Fläche gebraucht?
\begin{teile}
\teil Um einen runden Tisch wird eine Borte genäht. \hfill \kreuz{Umfang}\kreuz{Fläche}
\teil Ein rundes Beet wird mit Rasen eingesät. \hfill \kreuz{Umfang}\kreuz{Fläche}
\teil Eine runde Pizza wird ganz mit Tomatensoße bestrichen. \hfill \kreuz{Umfang}\kreuz{Fläche}
\teil Ein Rad rollt einmal ganz herum. Wie weit kommt es? \hfill \kreuz{Umfang}\kreuz{Fläche}
\teil Um einen runden Teich wird ein Zaun gebaut. \hfill \kreuz{Umfang}\kreuz{Fläche}
\end{teile}
\end{aufgabe}

\verfahren{Radius und Durchmesser}
\begin{aufgabe}{Ich kann den Durchmesser aus dem Radius berechnen und umgekehrt.}
Berechne den Durchmesser $d$ aus dem Radius $r$.
\begin{teilezwei}
\tz $r = 3\,\text{cm}$ \quad \feld[cm]{d} & \tz $r = 6\,\text{m}$ \quad \feld[m]{d} \\
\tz $r = 3{,}5\,\text{cm}$ \quad \feld[cm]{d} & \\
\anweisung{Berechne den Radius $r$ aus dem Durchmesser $d$.}
\tz $d = 10\,\text{cm}$ \quad \feld[cm]{r} & \tz $d = 14\,\text{m}$ \quad \feld[m]{r} \\
\tz $d = 9\,\text{cm}$ \quad \feld[cm]{r} & \tz $d = 9{,}6\,\text{dm}$ \quad \feld[dm]{r} \\
\tz $d = 1{,}2\,\text{m}$ \quad \feld[cm]{r} & \\
\end{teilezwei}
\end{aufgabe}

\verfahren{Kreise zeichnen}
\begin{aufgabe}{Ich kann Mittelpunkt, Radius und Durchmesser in einen Kreis einzeichnen (kennst du seit Klasse 5).}
Zeichne in den Kreis ein und beschrifte.
\begin{teile}
\teil einen Radius $r$ \quad (linker Kreis)
\teil einen Durchmesser $d$ \quad (mittlerer Kreis)
\teil zwei verschiedene Durchmesser \quad (rechter Kreis). Wo schneiden sie sich? \leerfeld
\end{teile}

\begin{kreis}[1.5]\mittelpunkt{M}\end{kreis}\hfill
\begin{kreis}[1.5]\mittelpunkt{M}\end{kreis}\hfill
\begin{kreis}[1.5]\mittelpunkt{M}\end{kreis}

\end{aufgabe}

\begin{aufgabe}{Ich kann einen Kreis mit dem Zirkel zeichnen.}
Zeichne um $M$ den Kreis mit dem Zirkel.
\begin{teile}
\teil Radius $r = 2\,\text{cm}$ \quad (links)
\teil Durchmesser $d = 3\,\text{cm}$ \quad (Mitte)
\teil Durchmesser $d = 4{,}6\,\text{cm}$ \quad (rechts)
\end{teile}

\zeichenplatz\hfill\zeichenplatz\hfill\zeichenplatz

\end{aufgabe}

\verfahren{Umfang und Durchmesser vergleichen}
\begin{aufgabe}{Ich kann Umfang und Durchmesser eines Kreises vergleichen.}
Die Tabelle zeigt gemessene Werte. Berechne $u : d$ und runde auf zwei Stellen nach dem Komma.

\sachtabelle{lccc}{Gegenstand & Durchmesser $d$ & Umfang $u$ & $u : d$}{Dose & $7{,}4\,\text{cm}$ & $23{,}2\,\text{cm}$ & \leerzelle \\ Teller & $24{,}0\,\text{cm}$ & $75{,}5\,\text{cm}$ & \leerzelle \\ Münze & $2{,}3\,\text{cm}$ & $7{,}2\,\text{cm}$ & \leerzelle \\ Eimer & $30{,}0\,\text{cm}$ & $94{,}0\,\text{cm}$ & \leerzelle}

\begin{teile}
\teil Ein Rad hat den Durchmesser $50\,\text{cm}$. Schätze mit der Tabelle seinen Umfang. \hfill \leerfeld[cm]
\end{teile}
\end{aufgabe}

\verfahren{Umfang berechnen}
\begin{aufgabe}{Ich kann den Umfang eines Kreises berechnen.}
Berechne den Umfang $u$ aus dem Durchmesser $d$. Nimm die $\pi$-Taste und runde auf eine Stelle nach dem Komma.
\begin{teilezwei}
\tz $d = 3\,\text{cm}$ \quad \feld[cm]{u} & \tz $d = 6\,\text{cm}$ \quad \feld[cm]{u} \\
\tz $d = 9\,\text{cm}$ \quad \feld[cm]{u} & \tz $d = 20\,\text{m}$ \quad \feld[m]{u} \\
\anweisung{Jetzt ist der Radius $r$ gegeben.}
\tz $r = 4\,\text{cm}$ \quad \feld[cm]{u} & \tz $r = 1{,}8\,\text{m}$ \quad \feld[m]{u} \\
\end{teilezwei}
\end{aufgabe}

\begin{aufgabe}{Ich kann aus dem Umfang den Durchmesser und den Radius berechnen.}
Berechne den Durchmesser $d$. Runde auf eine Stelle nach dem Komma.
\begin{teilezwei}
\tz $u = 47{,}1\,\text{cm}$ \quad \feld[cm]{d} & \tz $u = 88\,\text{m}$ \quad \feld[m]{d} \\
\anweisung{Berechne den Radius $r$.}
\tz $u = 37{,}7\,\text{cm}$ \quad \feld[cm]{r} & \tz $u = 100\,\text{cm}$ \quad \feld[cm]{r} \\
\end{teilezwei}
\end{aufgabe}

\begin{aufgabe}{Ich kann die Länge von Kreisrändern in Figuren und Körpern berechnen.}
Berechne die Länge. Runde auf eine Stelle nach dem Komma.
\begin{teile}
\teil gebogener Rand eines Halbkreises mit $d = 12\,\text{cm}$ \hfill \leerfeld[cm]
\teil gebogener Rand eines Halbkreises mit $r = 7\,\text{m}$ \hfill \leerfeld[m]
\teil ganzer Rand eines Halbkreises (Bogen und Durchmesser) mit $d = 12\,\text{cm}$ \hfill \leerfeld[cm]
\teil Eine zylinderförmige Dose hat den Radius $3{,}8\,\text{cm}$. Ein Etikett wird genau einmal um die Dose geklebt, ohne Überlappung. Berechne die Länge des Etiketts. (P10 2022 OS) \hfill \leerfeld[cm]
\end{teile}

\zylinder{1}{2}{r}{}

\end{aufgabe}

\begin{aufgabe}{Ich finde den Fehler in einer Umfangsrechnung.}
Finde den Fehler, benenne ihn und korrigiere die Rechnung.
\begin{teile}
\teil Kreis mit dem Radius $r = 7\,\text{cm}$: \rechnung{u &= \pi \cdot 7\,\text{cm} \\ &\approx 22{,}0\,\text{cm}}
\teil Kreis mit dem Radius $r = 6\,\text{cm}$: \rechnung{u &= \pi \cdot 6^2\,\text{cm} \\ &\approx 113{,}1\,\text{cm}}
\end{teile}
\end{aufgabe}

\begin{aufgabe}{Ich kann begründen, warum Umfang geteilt durch Durchmesser bei jedem Kreis gleich ist.}
Begründe ohne Rechnung.
\begin{teile}
\teil Ein großer Kreis hat einen doppelt so großen Durchmesser wie ein kleiner. Wie verändert sich der Umfang?
\teil Tom sagt: „Bei einem großen Kreis ist $u : d$ größer als bei einem kleinen.“ Hat er recht?
\end{teile}
\end{aufgabe}

\begin{aufgabe}{Ich kann ausrechnen, wie weit ein Rad rollt.}
Ein Fahrradreifen hat den Durchmesser $70\,\text{cm}$.
\begin{teile}
\teil Wie weit kommt das Rad bei einer Umdrehung? Runde auf eine Stelle. \hfill \leerfeld[cm]
\teil Wie weit kommt es bei 500 Umdrehungen? Gib das Ergebnis in m an, auf eine Stelle gerundet. \hfill \leerfeld[m]
\teil Wie viele Umdrehungen braucht das Rad für $5\,\text{km}$? \hfill \leerfeld
\end{teile}
\end{aufgabe}
